\subsection{Tổng quan}
Giả sử ta có n+1 số thực đôi một phân biệt $x_0<x_1< \dots< x_n$, tương ứng với chúng ta có $y_0, y_1, \dots< y_n$. Ta kì vọng tìm được hàm $f$ đi qua các điểm $(x_0, y_0), (x_1, y_1), \dots, (x_n, y_n)$, tức là 
\begin{equation}
	f(x_i)=y_i,\quad \text{với } i=0, 1, \dots, n.\\
	\label{eq: ràng buộc nội suy}
\end{equation}
\begin{figure}[htbp]
	\centering
\begin{tikzpicture}[
	x=1.25cm,
	y=1.15cm,
	every node/.style={font=\small}
	]
	
	%------------------------------------------------
	% Axes
	%------------------------------------------------
	\draw[thin] (0,-0.8) -- (0,4.5);
	\draw[thin] (-0.15,0) -- (8.15,0);
	
	%------------------------------------------------
	% Important coordinates
	%------------------------------------------------
	\coordinate (X0) at (1.0,0);
	\coordinate (P0) at (1.0,0.55);
	
	\coordinate (X1) at (2.4,0);
	\coordinate (P1) at (2.4,1.25);
	
	\coordinate (Xn) at (7.90,0);
	\coordinate (Pn) at (7.90,3.25);
	
	%------------------------------------------------
	% Main curve
	%------------------------------------------------
	\draw[thick]
	(-0.75,0.55)
	.. controls (-0.45,0.15) and (-0.20,0.00) .. (0,0)
	.. controls (0.40,0.00) and (0.70,0.28) .. (1.0,0.55)
	.. controls (1.45,0.95) and (1.95,1.25) .. (2.4,1.25)
	.. controls (3.05,1.25) and (3.50,0.72) .. (4.05,0);
	
	% dashed continuation of left curve
	\draw[dashed,thick]
	(4.05,0)
	.. controls (4.25,-0.28) and (4.45,-0.67) .. (4.65,-1.00);
	
	%------------------------------------------------
	% Right straight branch
	%------------------------------------------------
	
	% lower dashed continuation
	\draw[dashed,thick]
	(7.23,-1.05) -- (7.37,-0.55);
	
	% solid part
	\draw[thick]
	(7.40,-0.45) -- (7.90,3.25);
	
	% upper dashed continuation
	\draw[dashed,thick]
	(7.90,3.25) -- (8.03,4.20);
	
	%------------------------------------------------
	% Points
	%------------------------------------------------
	\fill (P0) circle (1.8pt);
	\fill (P1) circle (1.8pt);
	\fill (Pn) circle (2pt);
	
	%------------------------------------------------
	% Vertical dotted lines
	%------------------------------------------------
	\draw[dotted]
	(X0) -- (P0);
	
	\draw[dotted]
	(X1) -- (P1);
	
	\draw[dotted]
	(Xn) -- (Pn);
	
	%------------------------------------------------
	% Horizontal dotted lines
	%------------------------------------------------
	\draw[dotted]
	(0,0.55) -- (P0);
	
	\draw[dotted]
	(0,1.25) -- (P1);
	
	\draw[dotted]
	(0,3.25) -- (Pn);
	
	%------------------------------------------------
	% x-axis labels
	%------------------------------------------------
	\node[below] at (X0) {$x_0$};
	\node[below] at (X1) {$x_1$};
	
	% x_n is slightly to the left of the vertical dotted line
	\node[below left, xshift=3pt, yshift=-1pt] at (Xn) {$x_n$};
	
	\node[below] at (5.85,0) {$\cdots$};
	
	%------------------------------------------------
	% f labels
	%------------------------------------------------
	\node[left] at (0,0.55) {$y_0$};
	\node[left] at (0,1.25) {$y_1$};
	
	\node[left] at (0,2.00) {$\vdots$};
	
	\node[left] at (0,3.25) {$y_n$};
	
	%------------------------------------------------
	% label p'_l
	%------------------------------------------------
	\node[above left, xshift=4pt, yshift=3pt]
	at (8.00,4.05) {$f(x)$};
\end{tikzpicture}
	\caption{Các $(x_0, y_0), (x_1, y_1), \dots, (x_n, y_n)$ trên mặt phẳng cùng với hàm nội suy $f$}
	\label{fig:visulize_noi_suy}
\end{figure}

\textbf{Quy ước tên gọi: }
\begin{itemize}
	\item Các điểm $x_0, x_1, \dots, x_n$ ta gọi là điểm dữ liệu(hoặc mốc nội suy), các giá trị tương ứng $y_0, y_1, \dots, y_n$ ta gọi là các dữ liệu.
	\item Tập $S=\{(x_i, y_i): i=0, 1, \dots n \}$ được gọi là tập các điểm nội suy.
	\item  Hàm $f(x)$ được gọi là hàm nội suy, trong trường hợp $f(x)$ là đa thức theo biến $x$, ta gọi $f(x)$ là đa thức nội suy.
\end{itemize}
\begin{menhde}{}{}
Với $\mathbf{R}_n[x]$ là tập các đa thức hệ số thực có bậc không quá $n$, ta có:

$P_n(x)=a_n x^n+ a_{n-1}x^{n-1}+\dots a_1 x+a_0\in \mathbf{R}_n[x]$ và thoả \eqref{eq: ràng buộc nội suy} khi và chỉ khi hệ phương trình sau đươc thoả mãn
\begin{equation}
	\begin{bmatrix}
		1 & x_0 & x_0^2 & \cdots & x_0^n \\
		1 & x_1 & x_1^2 & \cdots & x_1^n \\
		\vdots & \vdots & \vdots & \ddots & \vdots \\
		1 & x_n & x_n^2 & \cdots & x_n^n 
	\end{bmatrix}	
\begin{bmatrix} 
a_0\\ a_1 \\ \vdots \\ a_n
\end{bmatrix}=
\begin{bmatrix}
y_0\\ y_1 \\ \vdots \\ y_n
\end{bmatrix}
\end{equation}
\end{menhde}

Rõ ràng rằng, phép nội suy nói chung và trường hợp nội suy bằng đa thức nói riêng, sai số là điều không thể tránh. Ta có định lý sau.
\begin{dinhly}{CT sai số}{ct-sai-so}
Cho $f\in C^{n+1}([a,b])$ và $P_n$ là đa thức nội suy của $f$ tại các nút đôi một phân biệt $x_0,x_1,\dots,x_n\in[a,b]$. Khi đó, với mỗi $x\in[a,b]$, tồn tại $\xi=\xi(x)\in I_x=(\min(x_0, x_1,\dots, x_n), \max(x_0, x_1, \dots, x_n)) \subseteq[a,b]$ sao cho
\begin{equation}
	f(x)-P_n(x)=\frac{f^{(n+1)}(\xi_x)}{(n+1)!}\prod_{i=0}^{n}(x-x_i).
	\label{eq:sai-so-noi-suy}
\end{equation}
\end{dinhly}
\begin{chungminh}
	\textbf{Trường hợp 1:} $x$ trùng với một trong các mốc nội suy. (Định lý hiển nhiên đúng)
	
	\noindent \textbf{Trường hợp 2:} $x \neq x_i, \quad \forall i=0, 1, \dots, n$.
	
	Cố định giá trị $x\in [a, b]$. Xét hàm 
	\[
	g(t)=f(t)-P_n(t)-K \omega _{n+1}(t), 
	\]
	với $K=\displaystyle \frac{f(x)-P_n(x)}{\omega_{n+1}(x)}, \omega_{n+1}(x) = \prod_{i=0}^n (x - x_i)$.
	
	Ta có:
	\begin{itemize}
		\item Với $\forall i =0, 1, \dots n$, $g(x_i)=f(x_i)-P_n(x_i)-K \omega_{n+1}(x_i)=0-K 0 =0$.
		\item $g(x)=f(x)-P_n(x)-\displaystyle \frac{f(x)-P_n(x)}{\omega_{n+1}(x)} \omega_{n+1}(x)=0$
	\end{itemize}
	Do đó, $g(t)$ có ít nhất $n+2$ nghiệm phân biệt trên đoạn đóng $\overline{I_x}$.
	
	Mặt khác, dễ dàng chứng minh được $g$ liên tục trên $\overline{I_x}$ và khả vi trên $I_x$. Áp dụng định lý Rolle ta được:
	\begin{itemize}
		\item $g'(t)=0$ có ít nhất $n+1$ nghiệm phân biệt trong $I_x$.
		\item $g''(t)=0$ có ít nhất $n$ nghiệm phân biệt trong $I_x$.
		\item ...
		\item Quy nạp ta được $g^{(n+1)}(t)=0$ có ít nhất $1$ nghiệm trong $I_x$.
	\end{itemize}
	
	Gọi $\xi$ là 1 nghiệm của $g^{(n+1)}(t)=0$. Khi đó, 
	\[
	g^{(n+1)}(\xi) = f^{(n+1)}(\xi) - K (n+1)! =0 \longrightarrow K = \frac{f^{(n+1)}(\xi)}{(n+1)!}.
	\]
	Như vậy, 
	\[
	\frac{f(x)-P_n(x)}{\omega_{n+1}(x)} = \frac{f^{(n+1)}(\xi)}{(n+1)!}.
	\]
	Nhân 2 vế với $\omega_{n+1}(x)$ ta được điều cần chứng minh.
\end{chungminh}


